\MFQLead{执行与做市控制答案}
\begin{enumerate}[leftmargin=*]
  \item 临时冲击 $\eta x_k$ 只进入第 $k$ 期价格；永久冲击 $\gamma X_{k-1}$ 随累计成交进入所有后续期。两者对“快做还是慢做”给出不同边际成本。
  \item 第一期价格 $100+0.2\times3=100.6$，成本 $1.8$；第二期价格 $100+0.1\times3+0.2\times2=100.7$，成本 $1.4$，合计 $3.2$，剩余 4 手另行报告。
  \item $V_K(0)=0$，$V_K(R>0)=+\infty$ 或显式终端罚函数。否则算法会通过不交易获得虚假零成本。
  \item 枚举每期 $0,1,2,3$，筛选总和为 6；按同一目标重算，最优 $(3,2,1)$、成本 $19.2$，与递归 DP 相同。
  \item $\kappa$ 或库存风险越高，多头状态下报价中心下降越多，ask 更积极、bid 更保守；但成交强度响应若未建模，不能断言期末库存必然下降。
  \item 若整条报价序列使用事后库存，则未来成交影响过去报价。正确循环必须是读取当前状态、报价、观察成交、更新状态，再进入下一步。
  \item 修复包至少含决策时到达价、方向、计划与实际量、逐笔成交、费用、未完成量、停止原因、基点与现金两种口径、对照基准和时间戳。
  \item 压力矩阵交叉扫描完成率分布、临时/永久冲击、期限和损失阈值；所有控制器复用相同随机样本，并报告完成率、IS、尾部和触停概率。
\end{enumerate}

事先确定 oracle：实际成交 5、剩余 4、partial-fill IS 为 $3.2$；DP 与完整枚举均给出 $(3,2,1)$ 和目标 $19.2$；bid 成交后下一轮报价为 $99.4/100.4$。

\section*{逐题推导与核对}
\begin{enumerate}[leftmargin=*]
  \item 设第 $k$ 期交易量为 $x_k$，累计库存为 $X_k=\sum_{j\le k}x_j$。成交价可写为 $m_k+\eta x_k+\gamma X_{k-1}$：$\eta x_k$ 只影响当前成交，$\gamma X_{k-1}$ 代表过去交易留下的永久位移并进入未来中间价。于是临时冲击鼓励分散单期交易，永久冲击还会改变后续每一期的基准，不能把两项合并成一个固定费率。
  \item 日程 $(3,2,0)$ 下，第一期到达价为 $100$，临时冲击使成交价 $100.6$，3 手的 IS 为 $3(100.6-100)=1.8$。累计成交 3 手带来下一期永久项 $0.1\times3$，第二期成交价为 $100+0.3+0.4=100.7$，2 手贡献 $1.4$，合计 $3.2$；剩余 1 手必须在期末报告为未完成，而不是把它按零价格计入。
  \item 令 $V_k(t,r)$ 表示时点 $t$ 还有库存 $r$ 时的最小未来成本。终端条件是 $V_K(0)=0$，但 $V_K(r>0)$ 应为无穷或明确的强制平仓罚函数；否则动态规划会选择永远不交易，把剩余库存“免费消失”。递推为
  \[
  V_k(t,r)=\min_{0\le x\le r}\{\text{期内成本}(x)+\mathbb E[V_{k+1}(t+1,r-x)]\}.
  \]
  \item 枚举六手三期时，每期量在 $0,\ldots,6$ 中取值并保留 $x_1+x_2+x_3=6$ 的组合。对每个组合按同一冲击和终端罚函数算目标，再取最小值；结果 $(3,2,1)$、目标 $19.2$ 与递归 DP 相同。逐项比较不仅检验最优值，也检验状态转移和边界条件。
  \item 提高库存惩罚 $\kappa$ 后，值函数中 $\kappa r^2$ 的边际梯度变大；多头时控制器更愿意压低报价中心、提高卖出强度，多头 bid 更保守。必须在相同随机完成率下比较报价、成交率、期末库存和尾部损失；只观察报价中心不能推出库存一定下降。
  \item 正确控制循环是：读取当前库存和历史，生成当前报价，观察当期成交，再更新库存与现金，进入下一期。若先生成整条报价序列再把最终库存倒填进去，报价使用了未来成交信息；负例可把未来一笔大成交注入前一期，检查输出是否改变。
  \item 修复执行报告时，逐笔保存到达价、方向、计划量、实际量、费用、未完成量、停止原因、时间戳和现金/基点两套口径。平均成交价只能回答已成交部分的价格，不能回答未完成风险或停机原因；因此还需对照 VWAP、arrival price 和强制平仓基准。
  \item 压力矩阵交叉扫描完成率、临时/永久冲击、剩余期限、库存惩罚和最大损失阈值。所有控制器使用同一随机完成率序列，输出完成率、IS 均值、P95/P99、触停概率和期末库存；这样比较的是策略而不是随机数。
\end{enumerate}
